%% ma: L10-B13-0078
%% lop: 10
%% chuong: 5
%% bai: 13
%% dang: 10.13.10
%% loai: TN4
%% muc: TH
%% dap_an: "A"
%% nguon: "chuyen-de-cac-so-dac-trung-cua-mau-so-lieu-khong-ghep-nhom-toan-10-chuong-trinh-moi.pdf (GV. Phan Nhật Linh) (P03), Bài 02 Dạng 2 Phần I Câu 3"
%% kien_thuc: "Lớp 10 Bài 13 (trung vị, tứ phân vị)"
%% trang_thai: chinh-thuc
%% ghi_chu: "Tự duyệt theo yêu cầu giáo viên 22/10/2026."
\begin{ex}
Chỉ số IQ của một nhóm học sinh là:
\begin{center}
\begin{tabular}{|c|c|c|c|c|c|c|c|c|c|}\hline
$60$ & $78$ & $80$ & $64$ & $70$ & $76$ & $80$ & $74$ & $86$ & $90$\\\hline
\end{tabular}
\end{center}
Các tứ phân vị của mẫu số liệu là
\choice{\True $Q_1=70;Q_2=77;Q_3=80$}{$Q_1=72;Q_2=78;Q_3=80$}{$Q_1=70;Q_2=76;Q_3=80$}{$Q_1=70;Q_2=75;Q_3=80$}
\loigiai{Sắp xếp: $60$; $64$; $70$; $74$; $76$; $78$; $80$; $80$; $86$; $90$.

$Q_2=\dfrac{76+78}2=77$; nửa dưới $60$; $64$; $70$; $74$; $76$ nên $Q_1=70$; nửa trên $78$; $80$; $80$; $86$; $90$ nên $Q_3=80$. Chọn A.}
\end{ex}
